\chapter{Valuative analytification of complex singularities}
\index{analytification!complex singularity}\index{semivaluation!multiplier-ideal formula}
\label{chap:valuative-analytification}

Multiplier ideals turn a complex psh singularity into a function on
semivaluations.
The construction is elementary at a fixed singularity, but two points need
care.
First, an ideal can vanish identically on the support of a semivaluation, so
the correct value is $+\infty$ before changing sign.
Second, the analytification of a test curve is initially defined by a
Legendre supremum on divisorial valuations; proving that it is a
Boucksom--Jonsson psh metric requires algebraization and uniform global
generation.

The chapter develops those steps in the order used in
\cite{DXZ23}.
The proof uses the positive-mass recovery of
\Cref{thm:tna-positive-volume}, the relative ray correspondence of
\Cref{thm:rays-rwn}, and the BJ closure and divisorial determination
of \Cref{thm:bj-basic-closure}. The published comparison section assumes
a big line-bundle class. Here the fixed-potential generation argument also
works in a pseudoeffective class, and the final perturbation step treats
zero mass explicitly.
Throughout, $X$ is a connected projective manifold over $\mathbf C$.
It has dimension $n$, $L$ is a pseudo-effective
$\mathbf Q$-line bundle, $h$ is a smooth reference metric, and
$\theta=c_1(L,h)$.

The sign ledger is fixed as follows.  Ideal values are nonnegative,
$\log|\mathfrak a|(v)=-v(\mathfrak a)$, and a complex singularity is
sent to a nonpositive function.  We normalize
$\nu_{\ord_E}(\log|f|^2)=\ord_E(f)$, use
$t=-\log|z|^2$, and normalize the Gauss extension by
$\sigma(v)(z)=1$.  With these choices the Legendre term is $+\tau$
in \eqref{eq:va-test-curve-an}, and no factor of two occurs in the
divisorial formula.

\section{Ideal values and the multiplier-ideal formula}
\label{sec:va-ideal-values}

A point $v\in X^{\An}$ is a semivaluation on an irreducible subvariety
$Y\subseteq X$, called its support.
Because $X$ is proper, it has a center on $X$.
For a coherent ideal $\mathfrak a$, define
\[
 v(\mathfrak a)
 :=\min\{v(f):f\text{ is a local generator of }\mathfrak a
                    \text{ at }c(v)\}\in[0,+\infty].
\tag{56.1}\label{eq:va-ideal-value}
\]
If $g_j=\sum_i a_{ij}f_i$ is another local generating family, then
$v(a_{ij})\geq0$ and the valuation inequality gives
$v(g_j)\geq\min_i v(f_i)$. Reversing the families proves independence.
These are the tropical ideal rules of
\Cref{eq:ideal-tropical-identities}: products add values, sums take minima,
and ideal inclusion reverses their order. Local units have value zero.
The minimum equals $+\infty$ exactly when the restriction of $\mathfrak a$ to
the support of $v$ is the zero ideal near the center.

\begin{definition}[Analytification of a psh singularity]
\label{def:va-psh-analytification}
For $\varphi\in\PSH(X,\theta)$, set
\[
 \varphi^{\An}(v)
 :=-v(\varphi)
 :=-\lim_{k\to\infty}\frac1k
          v\!\left(\cI(k\varphi)\right)
 \in[-\infty,0].
\tag{56.2}\label{eq:va-multiplier-limit}
\]
\end{definition}

The formula contains no background form. It therefore also defines
$\varphi^{\An}$ for every $\varphi\in\QPSH(X)$: choose a sufficiently
large positive multiple of an ample curvature form to make $\varphi$
psh, and use the same multiplier ideals.


\begin{proposition}\label{prop:va-fekete}
The limit in \eqref{eq:va-multiplier-limit} exists for every semivaluation,
including semivaluations of proper support, and
\[
 \varphi^{\An}(v)
   =\inf_{j\geq0}
      \left[-2^{-j}v\!\left(\cI(2^j\varphi)\right)\right].
\tag{56.3}\label{eq:va-dyadic}
\]
\end{proposition}

\begin{proof}
Put $a_0=0$ and $a_k=v(\cI(k\varphi))$ for $k\geq1$.
Subadditivity of multiplier ideals and the ideal product rule give
\[
 \cI((k+\ell)\varphi)\subseteq\cI(k\varphi)\cI(\ell\varphi),
 \qquad a_{k+\ell}\geq a_k+a_\ell.
\]
All values lie in $[0,+\infty]$. Fix $k\geq1$ and write
$N=qk+r$, $0\leq r<k$. Then
\[
                  a_N\geq qa_k+a_r\geq qa_k.
\]
If $a_k=+\infty$, this forces $a_N=+\infty$ for every $N\geq k$.
Otherwise, dividing by $N$ and letting $N\to\infty$ gives
$\liminf_N N^{-1}a_N\geq k^{-1}a_k$. Taking the supremum over $k$
and using the reverse limsup bound proves
\[
           \lim_{N\to\infty}N^{-1}a_N
           =\sup_{k\geq1}k^{-1}a_k.
\]
The same argument with $N=2^j$ proves that the dyadic subsequence
has this limit. Its terms $2^{-j}a_{2^j}$ are increasing by
$a_{2k}\geq2a_k$. Changing signs proves \eqref{eq:va-dyadic},
including the value $-\infty$ on proper supports.
\end{proof}

\begin{proposition}\label{prop:va-divisorial-formula}
For a prime divisor $E$ on a modification of $X$ and $t>0$,
\[
             \varphi^{\An}(t\ord_E)=-t\,\nu_E(\varphi),
\tag{56.4}\label{eq:va-divisorial}
\]
where $\nu_E(\varphi)$ is the generic Lelong number.
\end{proposition}

\begin{proof}
The asymptotic multiplier-ideal input is the divisorial Bergman
approximation estimate, recorded in the proof of
\cite[Theorem~5.1]{BFJ08}. In the squared-norm convention of this book,
if $A_X(E)=1+\ord_E(K_{Y/X})$ on a smooth model carrying $E$, it gives
\[
 k\nu_E(\varphi)-A_X(E)
 \leq\ord_E\cI(k\varphi)\leq k\nu_E(\varphi).
\]
The lower bound comes from change of variables and local integrability
at the generic point of $E$; the upper bound comes from the Bergman
approximation being less singular than the original potential. Smooth
background terms change neither quantity. Dividing by $k$ proves
\[
 \lim_{k\to\infty}k^{-1}\ord_E\cI(k\varphi)=\nu_E(\varphi).
\]
Homogeneity of ideal values under $t\ord_E$ proves the formula.
Notice that this is an asymptotic equality; no claim is made that
$\ord_E\cI(k\varphi)=k\nu_E(\varphi)$ at a fixed $k$.
\end{proof}

\begin{example}[A logarithmic pole on $\mathbf P^1$]
\label{ex:va-log-pole}
Let $X=\mathbf P^1$, let $L=\mathcal O(1)$ carry a smooth positive
metric $h$ of curvature $\theta$, and choose a section $s$ vanishing
simply at $p$.  For $0<\lambda\leq1$, set
\[
                       \varphi=\lambda\log|s|_h^2,
\]
with the Poincar\'e--Lelong normalization
$\nu_p(\log|s|_h^2)=1$.  Then
\[
              \theta+\ddc\varphi=(1-\lambda)\theta+\lambda[p]\geq0,
\]
so $\varphi\in\PSH(X,\theta)$.

If $z$ is a local coordinate at $p$, the one-variable multiplier-ideal
calculation \eqref{eq:complex-multiplier-example} gives
\[
                 \cI(k\varphi)_p=(z^{\lfloor k\lambda\rfloor}).
\]
For $v_c=c\,\ord_p$, with $c>0$, it follows that
\[
 \varphi^{\An}(v_c)
   =-\lim_{k\to\infty}\frac{c\lfloor k\lambda\rfloor}{k}
   =-c\lambda,
\]
in agreement with \Cref{prop:va-divisorial-formula}.  At the trivial
valuation the value is zero.  At a semivaluation supported on the point
$p$, the ideal restricts to zero once
$\lfloor k\lambda\rfloor>0$; its ideal value is then $+\infty$, and
$\varphi^{\An}=-\infty$.  Thus this example also exhibits the
proper-support convention in \eqref{eq:va-multiplier-limit}.
\end{example}

The formula immediately implies
\[
 \varphi\preceq_{\cI}\psi
       \quad\Longrightarrow\quad
 \varphi^{\An}\leq\psi^{\An}.
\tag{56.5}\label{eq:va-i-order}
\]
Conversely, comparison on all divisorial valuations recovers the
$\cI$-order by the valuative criterion for multiplier ideals.

\par\medskip\noindent\begin{minipage}{\linewidth}
\begin{proposition}
\label{prop:va-i-equivalence}
For $\varphi,\psi\in\QPSH(X)$, the following are equivalent:
\begin{enumerate}
\item $\cI(m\varphi)=\cI(m\psi)$ for every real $m>0$;
\item $\nu_E(\varphi)=\nu_E(\psi)$ for every prime divisor $E$ over
      $X$;
\item $\varphi^{\An}=\psi^{\An}$ on all divisorial semivaluations;
\item $\varphi^{\An}=\psi^{\An}$ on all of $X^{\An}$.
\end{enumerate}
Thus analytification records $\mathcal I$-equivalence.  It does not, for
arbitrary singularities, identify $P$-equivalence classes.
\end{proposition}
\end{minipage}\par\medskip

\begin{proof}
The first condition implies the fourth by
\eqref{eq:va-multiplier-limit}, and the fourth implies the third.
The third is equivalent to the second by
\Cref{prop:va-divisorial-formula}.  Conversely, equality of all divisorial
Lelong numbers gives equality of the upper-regularized multiplier ideals by
the valuative characterization \cite[Theorem~A]{BFJ08}.  Strong openness
\cite{GZ15} identifies those upper-regularized ideals with
$\cI(m\varphi)$ and $\cI(m\psi)$, proving the first condition.  The
last sentence is witnessed by the polar Cantor example in
\Cref{prop:ok-i-invariant-obstruction}: its potential is ideal-equivalent
to zero, but its raw mass is zero rather than one. The $P$-envelope
preserves raw mass, by \eqref{eq:complex-p-projection}, so their
$P$-classes differ. On the positive-mass ideal-good domain the two envelope
constructions agree, by \cite[Theorem~7.1.1]{XiaBook}.
\end{proof}

\section{Gauss extension and complexification}
\label{sec:va-gauss}

Let $z$ be the coordinate on $\mathbf A^1$; the letter $t$ remains
reserved for geodesic time.

\begin{definition}[Gauss extension]
\label{def:va-gauss-extension}
For a semivaluation $v$ on $X$, its Gauss extension is
\[
 \sigma(v)\!\left(\sum_i f_iz^i\right)
       :=\min_i\{v(f_i)+i\}.
\tag{56.6}\label{eq:va-gauss}
\]
\end{definition}

The kernel in \eqref{eq:va-gauss} consists precisely of polynomials
whose coefficients vanish on the support $Y$ of $v$. Thus the support
of $\sigma(v)$ is $Y\times\mathbf A^1$, and its center is
$c(v)\times\{0\}$. Zero coefficients have value $+\infty$.
After passing to the function field of $Y$, multiplicativity follows
by taking initial forms: the associated graded ring of a valued field
is a domain, and adjoining the initial form of $z$ preserves this
property. The ultrametric inequality follows directly from the minimum
formula. Hence $\sigma(v)$ is a semivaluation, including when $Y\ne X$.

The divisorial assertion uses the rational convention of
\Cref{def:ch32-divisorial}. If $v=(a/b)\ord_E$, where $a,b$ are coprime
positive integers, take a local equation $u$ of $E$ near its generic
point. The Gauss extension is monomial in $(u,z)$ with weights
$(a/b,1)$. The exceptional prime $D$ of the normalized weighted blowup
with weights $(a,b)$ therefore gives
\[
                   \sigma(v)=b^{-1}\ord_D.
\tag{56.6a}\label{eq:va-gauss-weighted-divisor}
\]
For the trivial valuation, the Gauss divisor is $X\times\{0\}$.
An irrational real multiple of $\ord_E$ still satisfies the ideal-value
formula, but its Gauss extension has rational rank two and is not
divisorial. This distinguishes \eqref{eq:va-divisorial}, which holds
for every real $t>0$, from the divisorial argument below.
Our normalization is
\[
                         \sigma(v)(z)=1,
 \qquad \nu_{\sigma(v)}(\log|z|^2)=1
\]
in the divisorial case. The corresponding geodesic time remains
$t=-\log|z|^2$.

Let $\Gamma=(\Gamma_\tau)$ be a transcendental non-Archimedean metric in
the class $\{\theta\}$.
It is represented, after adding an arbitrary K\"ahler form $\omega$, by an
$\cI$-model test curve
$\Gamma^{\theta+\omega}$.

\begin{definition}[Analytification of a test curve]
\label{def:va-test-curve-analytification}
For $v\in X^{\divis}$, define
\[
 \Gamma^{\An}(v)
 :=\sup_{\tau\leq\Gamma_{\max}}
       \left\{
       \bigl(\Gamma_\tau^{\theta+\omega}\bigr)^{\An}(v)+\tau
       \right\}.
\tag{56.7}\label{eq:va-test-curve-an}
\]
The convention is
$(-\infty)^{\An}(v)=-\infty$.
\end{definition}

The supremum in \eqref{eq:va-test-curve-an} can equally be taken over
$\tau<\Gamma_{\max}$, as in \cite[after (17)]{DXZ23}.
Indeed, put $c=\Gamma_{\max}$. If the terminal level is nontrivial,
monotonicity gives, for $s<c$,
\[
 \bigl(\Gamma_c^{\theta+\omega}\bigr)^{\An}(v)+c
 \leq\bigl(\Gamma_s^{\theta+\omega}\bigr)^{\An}(v)+s+(c-s).
\]
Letting $s\uparrow c$ shows that the terminal summand does not increase
the supremum. If the terminal level is $-\infty$, the assertion is
immediate. Thus the closed endpoint convention of
\Cref{eq:tna-component-compatibility} is compatible with this formula.

On $X^{\divis}$ this supremum is real: each nontrivial level has
finite generic divisorial Lelong numbers, at least one level is nontrivial,
and every summand is at most $\Gamma_{\max}$. At $v_{\triv}$ its value
is $\Gamma_{\max}$, since every nonzero multiplier ideal has value zero.
For a parameter translation $\Gamma'_\tau=\Gamma_{\tau-b}$,
\[
 (\Gamma')^{\An}=\Gamma^{\An}+b,\qquad
 \Gamma'_{\max}=\Gamma_{\max}+b.
\tag{56.7a}\label{eq:va-parameter-translation}
\]
This follows by reindexing the supremum and remains valid when a level
is $-\infty$.

Nothing in the notation $\Gamma^{\An}$ records $\omega$, so the
definition is only legitimate once the right side is known to be
independent of it.  That is the content of the next lemma, and it is the
reason the whole chapter may work with one arbitrary perturbation.

\begin{lemma}[The analytification does not see the perturbation]
\label{lem:va-omega-independence}
The right side of \eqref{eq:va-test-curve-an} is independent of the choice
of K\"ahler form $\omega$.
\Cref{def:va-test-curve-analytification} is therefore well posed.
\end{lemma}

\begin{proof}
Let $\omega\geq\omega'$ be K\"ahler and fix
$\tau<\Gamma_{\max}$. Compatibility
\eqref{eq:tna-component-compatibility} of the inverse system says that the
component at the larger perturbation is the $\cI$-projection of the
component at the smaller one,
\[
 \Gamma_\tau^{\theta+\omega}
 =P_{\theta+\omega}\bigl[\Gamma_\tau^{\theta+\omega'}\bigr]_{\cI}.
\]
Passing to an $\cI$-model representative is by construction an
$\cI$-equivalence, so
$\cI\bigl(k\Gamma_\tau^{\theta+\omega}\bigr)
  =\cI\bigl(k\Gamma_\tau^{\theta+\omega'}\bigr)$
for every $k\geq1$.  The implication
$\text{(i)}\Rightarrow\text{(iii)}$ of
\Cref{prop:va-i-equivalence} then gives
\[
 \bigl(\Gamma_\tau^{\theta+\omega}\bigr)^{\An}
 =\bigl(\Gamma_\tau^{\theta+\omega'}\bigr)^{\An}
\]
on divisorial valuations, for every such $\tau$. Each summand of the
strict-interior supremum is therefore unchanged. Its index set is
unchanged as well, because the terminal parameter
$\Gamma_{\max}$ does not depend on the perturbation, by
\Cref{prop:tna-common-terminal}. The preceding endpoint argument then
shows that the two suprema in \eqref{eq:va-test-curve-an} agree.

Two arbitrary K\"ahler forms $\omega_1,\omega_2$ are compared through a
common lower bound: a sufficiently small positive multiple of a fixed
K\"ahler form lies below both, and the previous paragraph applies to each
comparison in turn.
\end{proof}

For the Gauss computation, take the positive-mass representative
in its own background class supplied by \Cref{thm:tna-positive-volume},
and denote that background by $\theta$.
We also use the same divisorial Legendre formula for a relative curve
in the sense of \cite[Section~2.2]{DXZ23}: its parameter profile is
concave, decreasing and upper semicontinuous, with finite cutoff and
nontrivial initial potential. Such a curve need not have $P$-fixed
levels. This broader source terminology will be needed for PL
concavification below; it is not the maximal curve of
\Cref{def:rays-test-curve}.
Translate its parameter so that $\Gamma_{\max}\leq0$.
The complexification of the dual ray is the usc regularization of
\[
 \Psi(x,z)
  =\sup_{\tau\leq\Gamma_{\max}}
       \{\Gamma_\tau(x)-\tau\log|z|^2\},
 \qquad (x,z)\in X\times\Delta^*.
\tag{56.8}\label{eq:va-complexification}
\]
Each summand extends as a $p_1^*\theta$-psh function across $z=0$,
because $-\tau\geq0$. Moreover, all summands are bounded above by
$\Gamma_{-\infty}\circ p_1$ on the unit disk. This common local upper
bound shows that $\Psi$ extends as a $p_1^*\theta$-psh function on
$X\times\Delta$. A nontrivial level prevents its collapse to
$-\infty$.

\begin{proposition}\label{prop:va-gauss-identity}
For a relative curve $\Gamma$ with positive-mass initial potential
and $\Gamma_{\max}\leq0$, and every $v\in X^{\divis}$,
\[
                       \Gamma^{\An}(v)=-\sigma(v)(\Psi).
\tag{56.9}\label{eq:va-gauss-identity}
\]
\end{proposition}

\begin{proof}
Work on a model carrying the Gauss divisor from
\eqref{eq:va-gauss-weighted-divisor}. For a fixed nontrivial level,
the normalized generic Lelong number of the corresponding summand is
\[
 \sigma(v)\bigl(\Gamma_\tau-\tau\log|z|^2\bigr)
                         =v(\Gamma_\tau)-\tau.
\]
For $v=v_{\triv}$, the first term is zero and the same formula holds
along $X\times\{0\}$.

Here is the envelope step in \cite[Proposition~3.1]{DXZ23}.
Choquet reduction selects an increasing sequence of finite maxima
whose usc limit is $\Psi$. Its pullbacks to the model converge in
$L^1_{\mathrm{loc}}$ to the pullback of $\Psi$: the identification
holds off the exceptional locus and then follows everywhere by usc
extension. Generic Lelong numbers of a finite maximum are the minima
of those of its terms. Upper semicontinuity of Lelong numbers under
$L^1_{\mathrm{loc}}$ convergence, applied at a generic point of the
divisor, consequently gives
\[
 \sigma(v)(\Psi)\geq
       \inf_{\tau\leq\Gamma_{\max}}
          \{v(\Gamma_\tau)-\tau\}.
\]
The opposite inequality follows because $\Psi$ dominates every
summand. Changing signs gives \eqref{eq:va-test-curve-an}.
This argument identifies the function on $X^{\divis}$; its extension
to points with proper support will be constructed from ideals below.
\end{proof}

\begin{example}[The constant endpoint sign check]
\label{ex:va-constant-endpoint}
Assume $\theta$ is K\"ahler and let $a\leq0$.  Take
$\Gamma_\tau=0$ for $\tau\leq a$ and $-\infty$ for $\tau>a$.
Then $\Gamma^{\An}(v)=a$ for every full-support semivaluation, since all
multiplier-ideal values are zero.  The dual ray is $u_t=at$, while its
complexification is
$\Psi=-a\log|z|^2$.  Hence
\[
                   -\sigma(v)(\Psi)=-(-a)=a,
\]
which checks simultaneously the minus sign in
\eqref{eq:va-gauss-identity}, the $+\tau$ in
\eqref{eq:va-test-curve-an}, and the normalization
$\sigma(v)(z)=1$.
\end{example}

\begin{example}
\label{ex:va-hyperplane-analytification}
Take $X=\mathbf P^n$, $n\geq1$, $L=\mathcal O(1)$, and the curve
of \Cref{ex:rays-hyperplane-test-curve}. Write
$a_i=|x_i|^2/|x|^2$, $H=\{x_0=0\}$,
$P_0=[1:0:\cdots:0]$, and $\psi=\log a_0$.
Recall that $\Gamma_\tau=0$ for $\tau\leq0$,
$\Gamma_\tau$ has singularity type $\tau\psi$ for $0<\tau<1$,
and $\Gamma_1=\psi$. The terminal parameter is $1$.
Rational positive levels have analytic singularities in our convention;
real positive levels are $\mathcal I$-good by the argument in that
example. Thus the divisorial Legendre formula gives
\[
 \Gamma^{\An}(v)=\max\{0,1-v(\mathcal I_H)\},
                         \qquad v\in X^{\divis}.
\]
The levels with $\tau\leq0$ contribute $0$, and the positive levels
contribute the supremum of $\tau(1-v(\mathcal I_H))$ for
$0<\tau\leq1$.

This function has an explicit Fubini--Study extension to all of
$X^{\An}$, using the sections $x_0,\ldots,x_n$ of $L$ with weights
$1,0,\ldots,0$:
\[
 \begin{split}
 F(v)&=\max\{1-v(\mathcal I_H),-v(\mathcal I_{P_0})\}\\
     &=\max\{0,1-v(\mathcal I_H)\}.
 \end{split}
\tag{56.9a}\label{eq:va-hyperplane-fs}
\]
The second equality uses
$\mathcal I_H+\mathcal I_{P_0}=\mathcal O_X$: at least one of the
two ideal values is zero. In particular it remains valid when an
ideal vanishes on the support and has value $+\infty$.
The uniqueness statement in \Cref{thm:va-positive-mass-metric}
identifies $F$ with the full analytification of $\Gamma$.

To apply \eqref{eq:va-gauss-identity}, first shift the terminal
parameter to zero. Put $\Gamma'_\tau=\Gamma_{\tau+1}$.
Its ray is $u'_t=u_t-t$, and its complexification is
\[
 \Phi(x,z)=\log\left(a_0+|z|^2\sum_{i=1}^n a_i\right).
\tag{56.9b}\label{eq:va-hyperplane-shifted-potential}
\]
In a local frame of $L$, the holomorphic generators inside this
squared norm are $x_0,zx_1,\ldots,zx_n$. They generate the ideal
$\mathcal I_Z$ of the smooth codimension-two submanifold
$Z=H\times\{0\}$. Near $Z$ it is the ideal $(u,z)$, where $u$
is a local equation of $H$; away from $Z$ it is the unit ideal.
Hence
\[
 \sigma(v)(\mathcal I_Z)
   =\min\{v(\mathcal I_H),1+v(\mathcal I_{P_0})\}
   =\min\{v(\mathcal I_H),1\}.
\]
For divisorial $v$, this is $\sigma(v)(\Phi)$, so
$-\sigma(v)(\Phi)=F(v)-1$, as required by the shift.
The calculation evaluates holomorphic ideals; the smooth functions
$a_i$ themselves are not arguments of the semivaluation.

The multiplier ideals and their approximants can also be computed
at every stage. Locally along $Z$,
$\Phi=\log(|u|^2+|z|^2)+O(1)$. For an integer $m\geq1$,
radial integration in the two transverse complex variables shows
that a monomial of total degree $d$ is integrable against
$\exp(-m\Phi)$ precisely when $d>m-2$. Thus
\[
 \begin{split}
 \cI(m\Phi)&=\mathcal I_Z^{m-1},\\
 -m^{-1}\sigma(v)\bigl(\cI(m\Phi)\bigr)
   &=-\bigl(1-m^{-1}\bigr)\min\{v(\mathcal I_H),1\}.
 \end{split}
\tag{56.9c}\label{eq:va-hyperplane-multiplier-stages}
\]
The dyadic stages decrease from $0$ to $F-1$ on every semivaluation.
For support contained in $H$, the minimum is still $1$, so the
limit is $-1$ and the translated value is $F=0$.
At the full-support trivial valuation $F=1=\Gamma_{\max}$.
This verifies both the central-fiber normalization and the
proper-support values directly.
\end{example}

\section{Piecewise-linear test curves}
\label{sec:va-pl}

Fix numbers
\[
                  \tau_0>\tau_1>\cdots>\tau_N
\]
and potentials $A_i\in\PSH(X,\theta)$.
Let $A_\tau$ interpolate affinely between
$(\tau_i,A_i)$, be constant for $\tau\leq\tau_N$, and equal
$-\infty$ for $\tau>\tau_0$.

\begin{definition}
\label{def:va-pl-analytification}
Define its valuative function on every semivaluation by
\[
 A^{\An}(v):=\max_{0\leq i\leq N}\{A_i^{\An}(v)+\tau_i\},
                         \qquad v\in X^{\An}.
\tag{56.10}\label{eq:va-pl-max}
\]
On divisorial valuations this equals
$\sup_{\tau\leq\tau_0}(A_\tau^{\An}(v)+\tau)$.
The same equality on all semivaluations follows from
\Cref{prop:va-affine-compatibility}, proved below without using
PL concavification.
\end{definition}

For the divisorial equality, generic Lelong numbers are additive under
positive linear combinations, so $A_\tau^{\An}(v)+\tau$ is affine
on each parameter interval. It reaches its maximum at an endpoint.
For $\tau\leq\tau_N$ the potential is constant and the added parameter
is maximized at $\tau_N$. Endpoint values are evaluated directly,
without multiplying $-\infty$ by zero. The finite maximum itself is
always defined on proper supports, including when every term is
$-\infty$.

The interpolated curve need not itself be concave.  Assume that its
$-\infty$ level has positive mass, as will always be the case when this
construction is used below.  Let $\widetilde A$ be its upper concave
envelope in the Legendre sense: define
\[
 \check A_t:=\sup_{\tau\in\mathbf R}(A_\tau+t\tau)
            =\max_{0\leq i\leq N}(A_i+t\tau_i),
 \qquad
 \widetilde A_\tau:=\inf_{t>0}(\check A_t-t\tau).
\]
The new initial potential is $\max_i A_i$, since
$\check A_t\to\max_i A_i$ in $L^1$ as $t\downarrow0$.
It dominates $A_N$, so mass monotonicity retains positivity. The
initial potential can therefore change during concavification.
The finite maximum $\check A_t$ is a subgeodesic, since each
$A_i+t\tau_i$ is one.  The inverse Legendre theorem
for subgeodesics, in the form used in
\cite[Lemma~3.2]{DXZ23}, therefore makes $\widetilde A$ a relative curve in the broader
source sense just specified. It does not make every level $P$-fixed:
concavification and model maximalization are different operations.
The former is the imported analytic input here. No $P$- or
$\mathcal I$-projection is part of this concavification.

\begin{proposition}\label{prop:va-pl-replacement}
On divisorial valuations,
\[
                         \widetilde A^{\An}=A^{\An}.
\tag{56.11}\label{eq:va-pl-replacement}
\]
\end{proposition}

\begin{proof}
Translate the parameters so that $\tau_0\leq0$, and let $\Psi$ be the
$S^1$-invariant potential corresponding to the subgeodesic
$\check A_t$.  The same generic-Lelong-number computation as in
\Cref{prop:va-gauss-identity} gives
\[
 A^{\An}(v)=-\sigma(v)(\Psi)
\]
for divisorial $v$.  The curve $\widetilde A$ has the same dual
subgeodesic $\check A$, by scalar Legendre biconjugacy.  Applying
\Cref{prop:va-gauss-identity} to $\widetilde A$ therefore gives
$\widetilde A^{\An}(v)=-\sigma(v)(\Psi)$, proving
\eqref{eq:va-pl-replacement}.
\end{proof}

\begin{example}[Two poles and a PL curve]
\label{ex:va-pl-two-poles}
On $\mathbf P^1$ take distinct points $p,q$, the Fubini--Study
background $\theta$, and sections $s_p,s_q$ of $\mathcal O(1)$
vanishing there. For $0<\lambda<1$, put
\[
 A_0=\lambda\log|s_p|^2,\qquad
 A_1=\lambda\log|s_q|^2,\qquad
 \tau_0=1,\quad\tau_1=0,
\]
where the norms use the Fubini--Study metric. Both currents have raw
mass $1-\lambda>0$, by the pole calculation above. Their PL
valuative function is
\[
 A^{\An}(v)=\max\{1-\lambda v(\mathcal I_p),
                         -\lambda v(\mathcal I_q)\}.
\tag{56.11a}\label{eq:va-pl-two-poles}
\]
The pole ideal formula proves each term on every semivaluation,
including infinite ideal values. Since
$\mathcal I_p+\mathcal I_q=\mathcal O_X$, at least one ideal value
is zero. Thus $0\leq A^{\An}\leq1$. It takes value $1$ at
$v_{\triv}$, value $0$ at the point-support semivaluation of $p$,
and value $1$ at that of $q$. For $c\ord_p$ its value is
$\max(1-\lambda c,0)$.

The initial potential after concavification is
$\max(A_0,A_1)$. It is bounded because the poles are disjoint, so
its raw mass is one; the original left-tail potential $A_1$ had mass
$1-\lambda$. This explicitly shows why the concave envelope need
not retain the original initial potential. The divisorial function
is nevertheless preserved by \eqref{eq:va-pl-replacement}, and its
psh extension is the finite maximum in \Cref{cor:va-pl-bj-psh}.
\end{example}

\section{From multiplier ideals to Fubini--Study functions}
\label{sec:va-fs}

We now prove that the valuative functions above are psh metrics in the sense
of Boucksom--Jonsson.

\begin{theorem}\label{thm:va-single-potential-bj-psh}
For every $\varphi\in\PSH(X,\theta)$,
\[
                          \varphi^{\An}\in\PSH(L^{\An}).
\tag{56.12}\label{eq:va-single-psh}
\]
\end{theorem}

\begin{proof}
Choose a positive integer $r$ such that $rL$ is a line bundle.
Equation \eqref{eq:va-multiplier-limit}, evaluated along the subsequence
$rk$, gives $(r\varphi)^{\An}=r\varphi^{\An}$ on every
semivaluation. Thus homogeneity of the BJ cone reduces the assertion
to $(rL,r\theta,r\varphi)$;
we reset notation and assume that $L$ is integral.
Fix a very ample line bundle $H$, and choose $b$ so large that
$H^b\otimes K_X^{-1}\otimes H^{-q}$ is ample for
$1\leq q\leq n$. Nadel vanishing applied to the singular metric
on $L^k$ with weight $k\varphi$, together with each of these fixed
positive twists, gives
\[
 H^q\bigl(X,H^{b-q}\otimes L^k\otimes\cI(k\varphi)\bigr)=0
                       \quad(1\leq q\leq n).
\]
Castelnuovo--Mumford regularity therefore makes
\[
 H^b\otimes L^k\otimes\cI(k\varphi)
\tag{56.13}\label{eq:va-global-generation}
\]
globally generated for every $k\geq1$, with the same $b$.
This is the uniform-generation input of
\cite[Theorem~6.27]{DemAMAGAuthor2011} and \cite[Lemma~3.5]{DXZ23}.

Choose global generators $s_{k,j}$ and read them in a local frame
of $H^b\otimes L^k$ at the center of $v$. Changes of frame are
units and have value zero. Global generation then gives, on every
semivaluation,
\[
 v(\cI(k\varphi))=\min_jv(s_{k,j}).
\tag{56.14}\label{eq:va-generator-min}
\]
If the ideal vanishes on the support, both sides are $+\infty$.
Consequently
\[
 \phi_k(v):=-k^{-1}v(\cI(k\varphi))
       =k^{-1}\max_j\log|s_{k,j}|(v)
\tag{56.15}\label{eq:va-gf-stage}
\]
is a rational generically finite Fubini--Study function in the class
$L+(b/k)H$. By \eqref{eq:va-dyadic}, $\phi_{2^j}$ decreases
pointwise to $\varphi^{\An}$, while the auxiliary classes decrease
to $L$. The varying-class closedness clause of \Cref{thm:bj-basic-closure},
equivalently \cite[Theorem~4.5, property (PSH2)]{BJ22}, therefore gives
$\varphi^{\An}\in\PSH(L^{\An})$.
At the trivial valuation every nonzero multiplier ideal has value
zero, so the limit is not identically $-\infty$.
\end{proof}

\par\medskip\noindent\begin{minipage}{\linewidth}
\begin{proposition}[Affine compatibility]
\label{prop:va-affine-compatibility}
For $A,B\in\PSH(X,\theta)$ and $0<s<1$,
\[
 ((1-s)A+sB)^{\An}=(1-s)A^{\An}+sB^{\An}
                          \quad\text{on }X^{\An}.
\tag{56.15a}\label{eq:va-affine-compatibility}
\]
Consequently, the finite maximum in \eqref{eq:va-pl-max} equals its
Legendre supremum on all semivaluations.
\end{proposition}
\end{minipage}\par\medskip

\begin{proof}
Both sides are BJ psh in $L$: the left side by
\Cref{thm:va-single-potential-bj-psh}, the right side by the
positive-linear-combination clause of \Cref{thm:bj-basic-closure}.
Generic Lelong numbers are linear under positive combinations of
positive currents, so \eqref{eq:va-divisorial} identifies their
values on $X^{\divis}$. Divisorial determination in
\Cref{thm:bj-basic-closure} identifies them everywhere.
This proves additivity using psh extensions, rather than a false
finite-level multiplier-ideal equality.
Apply it on each open interpolation interval. An affine expression
with finite endpoint values has its supremum at an endpoint; if an
endpoint is $-\infty$, every interior combination involving it is
$-\infty$ and the other endpoint still controls the supremum.
At endpoints evaluate the specified potential directly. The constant
left tail is maximized at $\tau_N$. This proves the final claim,
including proper supports.
\end{proof}

\begin{corollary}\label{cor:va-pl-bj-psh}
The function $A^{\An}$ in \eqref{eq:va-pl-max} belongs to
$\PSH(L^{\An})$.
\end{corollary}

\begin{proof}
Each $A_i^{\An}$ is psh by
\Cref{thm:va-single-potential-bj-psh}.
Boucksom--Jonsson psh metrics are stable under translations and finite
maxima \cite[Theorem~4.7]{BJ22}.
\end{proof}

\section{Algebraizing the complexified multiplier ideals}
\label{sec:va-algebraization}

For general test curves, finite maxima are replaced by multiplier ideals on
$X\times\Delta$.
Two elementary algebraization facts isolate the geometry.

\begin{lemma}\label{lem:va-circle-ideal}
Let $R$ be a finitely generated $\mathbf C$-algebra and
$I\subset R[z]$ an ideal invariant under $z\mapsto\rho z$ for every
$\rho\in S^1$.
Then $I$ is $\mathbf C^*$-invariant and
\[
             I=I_0+I_1z+\cdots+I_m(z^m)
\tag{56.16}\label{eq:va-ideal-expansion}
\]
for ideals $I_0\subseteq I_1\subseteq\cdots\subseteq I_m$ of $R$.
\end{lemma}

\begin{proof}
If $a=\sum_{i=0}^Na_iz^i\in I$, evaluate its circle translates at
$N+1$ distinct points.
The Vandermonde matrix is invertible, hence each $a_iz^i$ belongs to
$I$.
Thus $I=\sum_{i\geq0}I_iz^i$ with $I_i\subseteq I_{i+1}$, because
multiplication by $z$ preserves $I$.
Noetherianity stabilizes this chain, giving
\eqref{eq:va-ideal-expansion}.
The homogeneous expansion also proves $\mathbf C^*$-invariance.
\end{proof}

\begin{lemma}\label{lem:va-gaga-ideal}
Let $\mathfrak A$ be a coherent analytic ideal on
$X\times\mathbf C$.
Assume
\[
 \mathfrak A|_{X\times\mathbf C^*}
       =p_1^*\mathfrak J|_{X\times\mathbf C^*}
\]
for an algebraic coherent ideal $\mathfrak J$ on $X$.
Then $\mathfrak A$ is algebraic.
\end{lemma}

\begin{proof}
Glue $\mathfrak A$ near $0$ to
$p_1^*\mathfrak J$ on
$X\times(\mathbf P^1\setminus\{0\})$.
This gives a coherent analytic ideal on the projective variety
$X\times\mathbf P^1$.
GAGA algebraizes it, and restriction to
$X\times\mathbf C$ algebraizes $\mathfrak A$.
\end{proof}

Let $\Psi$ be the complexification of the geodesic ray dual to
a positive-mass model test curve with $\Gamma_{\max}\leq0$.
Write $\phi=\Gamma_{-\infty}$ for its initial model potential.
We need an equality of ideals on a product, not just on its fibers:
\[
 \cI(m\Psi)|_{X\times\Delta^*}
       =p_1^*\cI(m\phi)|_{X\times\Delta^*}.
\tag{56.17}\label{eq:va-punctured-ideal}
\]
To prove it, fix $0<a<b$ and the annulus
$\Delta_{a,b}=\{e^{-b}<|z|^2<e^{-a}\}$.
The slices $u_a,u_b$ of the ray lie in
$\mathcal E(X,\theta;\phi)$ by \Cref{thm:rays-rwn}.
Relative rooftop stability and geodesic comparison, as used in
\Cref{thm:rays-relative-geodesic}, give
\[
 w:=P_\theta(u_a,u_b)\in\mathcal E(X,\theta;\phi),
 \qquad w\circ p_1\leq\Psi\leq\phi\circ p_1
                       \quad\hbox{on }X\times\Delta_{a,b}.
\]
The upper bound follows also from $\Gamma_\tau\leq\phi$ and
$\tau\leq0$. Relative full mass implies
$\cI(mw)=\cI(m\phi)$ for every $m$; this is the full-mass
multiplier-ideal input in the proof of
\cite[Proposition~3.10]{DXZ23}. It uses increasing rooftop
approximants and strong openness, and does not require bounded
equivalence of $w$ and $\phi$.
The local product formula for multiplier ideals gives
$\cI(mw\circ p_1)=p_1^*\cI(mw)$, and likewise for $\phi$;
see the same proof of \cite[Proposition~3.10]{DXZ23}.
The bounds thus yield
\[
 p_1^*\cI(mw)\subseteq\cI(m\Psi)
                      \subseteq p_1^*\cI(m\phi)
                      \quad\hbox{on }X\times\Delta_{a,b}.
\]
The outer ideals agree. Varying the annulus proves
\eqref{eq:va-punctured-ideal}.

Glue $\cI(m\Psi)$ to $p_1^*\cI(m\phi)$ on $X\times\mathbf C^*$.
The resulting ideal on $X\times\mathbf C$ is algebraic by
\Cref{lem:va-gaga-ideal} and retains its circle invariance.
Applying \Cref{lem:va-circle-ideal} on affine charts gives
\[
 \cI(m\Psi)=
 \mathfrak a_{m,0}+\mathfrak a_{m,1}z+\cdots
             +\mathfrak a_{m,N_m}(z^{N_m}),
\tag{56.18}\label{eq:va-multiplier-expansion}
\]
where the coefficient ideals are algebraic and increasing.
Here the left side denotes the extended ideal as well as its
restriction to the disk. Inverting $z$ and using stabilization shows
that $\mathfrak a_{m,N_m}=\cI(m\phi)$; this need not be
$\mathcal O_X$ when the initial model potential is unbounded.

\begin{lemma}\label{lem:va-family-generation}
Assume that $L$ is integral. There is an ample line bundle $T$,
independent of $m$, such that every
\[
                    T\otimes L^m\otimes\mathfrak a_{m,i}
\tag{56.19}\label{eq:va-coefficient-generated}
\]
is globally generated.
\end{lemma}

\begin{proof}
Positive mass of $\phi$ implies that $L$ is big. The uniform family
generation theorem \cite[Lemma~3.9]{DXZ23} therefore applies to
$\Psi\in\PSH(X\times\Delta,p_1^*\theta)$.
It supplies an ample $T$, independent of $m$, such that
$p_1^*(T\otimes L^m)\otimes\cI(m\Psi)$ is globally generated
on $X\times\Delta$. The analytic input is the relative vanishing
theorem \cite[Theorem~1.7]{Mat16}, with the fixed positive twists
in the uniform regularity argument. We use the family generation
statement in this precise form.

It remains to extract coefficient generators. A global section $S$
of this sheaf has a Taylor expansion in $z$. Its coefficient $S_\ell$
is a global section of $T\otimes L^m\otimes\mathfrak a_{m,\ell}$,
by the homogeneous expansion \eqref{eq:va-multiplier-expansion}.
Fix a local frame of $T\otimes L^m$ near $x$. For a local
section $f$ of $\mathfrak a_{m,i}$, express
$fz^i$ in the germs at $(x,0)$ of finitely many global generators
$S_r$, with holomorphic coefficients $h_r(x,z)$.
Taking the coefficient of $z^i$ expresses $f$ as an
$\mathcal O_{X,x}$-linear combination of $(S_r)_\ell$ with
$\ell\leq i$. Since
$\mathfrak a_{m,\ell}\subseteq\mathfrak a_{m,i}$, these are
global sections of the sheaf in \eqref{eq:va-coefficient-generated}.
They generate its stalk at $x$. Compactness of $X$ selects finitely
many such sections, and GAGA identifies them with algebraic sections.
The same $T$ works for all $m$ and $i$.
\end{proof}

\begin{theorem}\label{thm:va-positive-mass-metric}
If $\Gamma$ has positive mass, the function
\eqref{eq:va-test-curve-an} has a unique extension
\[
                           \Gamma^{\An}\in\PSH(L^{\An}).
\tag{56.20}\label{eq:va-positive-extension}
\]
\end{theorem}

\begin{proof}
Choose a positive integer $r$ with $rL$ integral. Homogeneity reduces
the assertion to the rescaled curve $r\Gamma$ in $rL$; divide the
resulting extension by $r$. We reset notation and assume that $L$
is a line bundle. Use \Cref{thm:tna-positive-volume} to take the positive-mass
model representative in the base class,
translate so that $\Gamma_{\max}\leq0$, and form $\Psi$.
Define on every semivaluation
\[
 \phi_m(v):=-\frac1m\sigma(v)(\cI(m\Psi))
 =-\frac1m\min_i\{v(\mathfrak a_{m,i})+i\}.
\tag{56.21}\label{eq:va-family-gf}
\]
Choose generators $s_{m,i,j}$ of
$T\otimes L^m\otimes\mathfrak a_{m,i}$ using
\Cref{lem:va-family-generation}. Then
\[
          \phi_m=\frac1m\max_{i,j}\{\log|s_{m,i,j}|-i\}.
\tag{56.21a}\label{eq:va-family-section-weights}
\]
Thus $\phi_m$ is a rational GF function in $L+m^{-1}T$, with
section weights $-i$. If every coefficient ideal vanishes on the
support of $v$, its value is $-\infty$; no support is discarded.

Subadditivity gives
$\cI(2m\Psi)\subseteq\cI(m\Psi)^2$.
The Gauss ideal value is therefore at least twice as large, and the
minus sign makes $\phi_{2m}\leq\phi_m$ on all semivaluations.
Identification of the divisorial limit uses the separate asymptotic
multiplier-ideal theorem obtained from Ohsawa--Takegoshi extension:
\[
 -\sigma(v)(\Psi)
 =\lim_{j\to\infty}
   -2^{-j}\sigma(v)\!\left(\cI(2^j\Psi)\right),
                           \qquad v\in X^{\divis}.
\]
This is the analytic input in \cite[Proposition~3.10]{DXZ23}.
By \Cref{prop:va-gauss-identity}, the limit is the prescribed
divisorial function. In particular, its value at $v_{\triv}$ is
$\Gamma_{\max}>-\infty$. The decreasing limit therefore does not
collapse. Since the auxiliary classes $L+2^{-j}T$ tend to $L$,
the BJ closedness axiom makes it a psh metric in $L$.
Finally undo the parameter translation.

For uniqueness, let $F$ and $G$ be two BJ psh extensions. The
divisorial domination theorem \cite[Theorem~4.22]{BJ22} applied
first to $(F,G)$ and then to $(G,F)$ gives $F\leq G$ and
$G\leq F$ on $X^{\An}$. Both functions are psh; equality with
an arbitrary usc extension would not follow from the one-sided
domination theorem.
\end{proof}

\section{The zero-mass safeguard}
\label{sec:va-zero-mass}

The complexification proof used positive mass.
It would be circular to assert that an arbitrary zero-mass metric has a
positive-mass ray.
Instead perturb the cohomology class.

\begin{theorem}\label{thm:va-all-masses}
For every transcendental non-Archimedean metric
$\Gamma$ in the pseudo-effective class $c_1(L)$, the divisorial function
\eqref{eq:va-test-curve-an} extends uniquely to
\[
                           \Gamma^{\An}\in\PSH(L^{\An}).
\tag{56.22}\label{eq:va-all-mass-extension}
\]
\end{theorem}

\begin{proof}
Choose an ample line bundle $H$ with K\"ahler form $\omega$.
For every $q\geq1$, the component
$\Gamma^{\theta+q^{-1}\omega}$ has positive mass by
\Cref{def:tna-pshna}. It is a curve in the rational big class
$c_1(L+q^{-1}H)$, so the preceding positive-mass construction applies.
\Cref{thm:va-positive-mass-metric} therefore gives
\[
 \Gamma^{\An}\in
       \PSH\bigl((L+q^{-1}H)^{\An}\bigr).
\tag{56.23}\label{eq:va-perturbed-classes}
\]
These extensions agree first on $X^{\divis}$ by
\Cref{lem:va-omega-independence}. For two choices of $q$, both
extensions are psh in the larger of their two background classes:
its difference from the smaller class is nef, and the constant zero
is psh in that difference by \Cref{lem:ch34-psh-nonempty}.
Adding it uses the sum rule in \Cref{thm:bj-basic-closure}.
Applying \cite[Theorem~4.22]{BJ22} in both directions shows that
they agree on $X^{\An}$, including points with proper support.
Thus there is one fixed extended-real function which is psh in every class
$L+q^{-1}H$.  The closedness axiom PSH2 of \cite{BJ22} lets
$q\to\infty$ and gives
\eqref{eq:va-all-mass-extension}.

At no point is a zero ideal removed.
If every coefficient ideal in
\eqref{eq:va-family-gf} vanishes on the support of $v$, its Gauss value is
$+\infty$ and the psh value is $-\infty$.
At the trivial full-support valuation,
$\Gamma^{\An}(v_{\triv})=\Gamma_{\max}>-\infty$, so the final metric is not
identically $-\infty$.
\end{proof}

\section{What the construction has proved}
\label{sec:va-conclusion}

The analytification map
\[
 \Gamma\longmapsto\Gamma^{\An}
\tag{56.24}\label{eq:va-an-map}
\]
is now defined on every transcendental non-Archimedean metric in a
pseudo-effective rational projective class.
Its values are extended-real functions on the whole trivially valued
analytification, not merely numerical data on divisorial points.
Perturbation independence proves compatibility with $\cI$-projection.
Equation \eqref{eq:va-parameter-translation} proves parameter-translation
compatibility on divisorial points and uniqueness extends it everywhere.
The finite PL maximum is given by \eqref{eq:va-pl-max}.

Surjectivity, order reflection, general monotone limits, operations, and
trace are not consequences of the construction alone.
They require the projective comparison theorem and are treated in the next
chapter.

\par\medskip\noindent\begin{minipage}{\linewidth}
\begin{remark}
\label{rem:va-source-boundary}
The divisorial Gauss identity is \cite[Proposition~3.1]{DXZ23}.
The algebraization and approximation follow
\cite[Lemmas~3.5 and~3.7--3.9, Proposition~3.10 and Corollary~3.11]{DXZ23}.
Uniform generation uses \cite[Theorem~6.27]{DemAMAGAuthor2011} in the fixed
case and \cite[Theorem~1.7]{Mat16} in the family case.
The PL step uses the unmaximalized relative-curve convention of
\cite[Section~2.2 and Lemma~3.2]{DXZ23}; it does not claim that concave
envelopes are model envelopes. Affine compatibility on proper supports
uses the fixed-potential theorem and divisorial determination.
The uniform-generation locator refers to the identified July 2011
author PDF, separately from the 2012 print edition.
The annulus argument retains relative full-mass rooftop stability
and full-mass multiplier-ideal equality as analytic inputs.
The target cone, decreasing GF closedness, and divisorial domination
are those of \cite[Theorems~4.5 and~4.22]{BJ22}.
This chapter constructs the extension and proves its uniqueness.
Surjectivity and compatibility with the Xia operations require the
comparison theorems of the next chapter.

\end{remark}
\end{minipage}\par\medskip
